% =========================================================
% Relatório — Edição de Atributos Faciais com Preservação de
% Identidade via CFG Composicional em Modelos de Difusão Latente
%
% Compilar com:  pdflatex relatorio_ldm_cfg.tex  (2x)
% Formato: artigo em duas colunas, estilo paper de ML (arXiv-like)
% =========================================================
\documentclass[10pt,twocolumn]{article}

\usepackage[utf8]{inputenc}
\usepackage[T1]{fontenc}
\usepackage[brazilian]{babel}
\usepackage{mathptmx}          % Times para texto e matemática
\usepackage{amsmath,amssymb}
\usepackage{booktabs}
\usepackage{graphicx}
\usepackage{microtype}
\usepackage[a4paper,margin=2cm,columnsep=0.7cm]{geometry}
\usepackage[colorlinks=true,linkcolor=black,citecolor=black,urlcolor=blue]{hyperref}
\usepackage{enumitem}
\setlist{leftmargin=*,itemsep=2pt,topsep=3pt}

% ---------------------------------------------------------
% Cabeçalho do paper
% ---------------------------------------------------------
\title{\LARGE\bfseries Edição de Atributos Faciais com Preservação de
Identidade via Classifier-Free Guidance Composicional em Modelos de
Difusão Latente}

\author{Lucas\\[2pt]
\normalsize\url{https://github.com/Lucas-023/diffusion-model}}

\date{}

\begin{document}

\twocolumn[
  \maketitle
  \begin{center}
  \begin{minipage}{0.85\textwidth}
    \small
    \begin{center}\textbf{Resumo}\end{center}
    Este trabalho investiga a edição de atributos faciais com preservação
    de identidade usando um Modelo de Difusão Latente (LDM) condicional
    treinado no CelebA. Dada a foto de uma pessoa, o sistema extrai um
    embedding de identidade (ArcFace), tokens de aparência (CLIP) e um
    vetor de 40 atributos binários; o usuário edita os atributos desejados
    e o modelo gera uma nova imagem que preserva a identidade original. A
    hipótese inicial --- combinar Classifier-Free Guidance (CFG) para
    identidade com Classifier Guidance (CG) para atributos (``Mixed
    Guidance'') --- mostrou-se estruturalmente frágil: a saturação da
    sigmoide no classificador multi-rótulo, a diluição do gradiente entre
    os 40 atributos e a janela temporal estreita em que um classificador
    sobre latentes ruidosos é informativo enfraquecem o sinal de edição.
    Substituímos o CG por \textbf{CFG composicional}~\cite{liu2022}, tratando
    identidade e atributos como condicionantes independentes com escalas de
    guidance separadas ($s_{\mathrm{id}}$, $s_{\mathrm{attr}}$). Duas
    contribuições adicionais foram necessárias na prática:
    (i)~\textbf{pareamento de referência por identidade} --- condicionar o
    treino em uma foto \emph{diferente} da mesma pessoa, impedindo que o
    ramo CLIP vaze a expressão da imagem-alvo e anule o token de atributo;
    e (ii)~um modo \textbf{split} com três ramos de guidance (ArcFace,
    CLIP, atributos), que permite zerar a contribuição de aparência da
    referência em edições estruturais (p.\,ex.\ \emph{Bald}) sem perder
    identidade. A avaliação combina FID periódico com amostragem DDIM,
    testes qualitativos de edição forçada (\emph{Smiling}=1) e um
    classificador de atributos CLIP ViT-L/14 ajustado no CelebA
    (F1\,$\approx$\,0{,}73) para verificação.

    \vspace{4pt}
    \textbf{Palavras-chave:} modelos de difusão latente, classifier-free
    guidance, edição de atributos faciais, preservação de identidade,
    CelebA, ArcFace.
  \end{minipage}
  \end{center}
  \vspace{18pt}
]

% =========================================================
\section{Introdução}
% =========================================================

Modelos de difusão tornaram-se o padrão para síntese de imagens de alta
qualidade~\cite{ho2020,rombach2022}, e sua variante latente (LDM)
\cite{rombach2022} reduz drasticamente o custo computacional ao operar no
espaço comprimido de um autoencoder. Um problema em aberto de interesse
prático é a \emph{edição condicional com preservação de identidade}: dada
a foto de uma pessoa, gerar a mesma pessoa com um subconjunto de
atributos faciais alterado --- fazê-la sorrir, adicionar óculos, remover
cabelo --- sem alterar quem ela é.

O pipeline proposto opera em cinco etapas: (1)~a imagem de entrada é
codificada no espaço latente por um VAE; (2)~a identidade é representada
por um embedding ArcFace~\cite{deng2019}; (3)~um classificador prediz os
40 atributos binários do CelebA a partir da imagem; (4)~o usuário
modifica um ou mais atributos preditos; (5)~o modelo de difusão gera a
nova imagem condicionada ao embedding de identidade e ao vetor de
atributos editado.

A questão de pesquisa original era se a combinação de Classifier-Free
Guidance~\cite{ho2022cfg} (para identidade) com Classifier
Guidance~\cite{dhariwal2021} (para atributos) melhora a controlabilidade
preservando identidade. Os experimentos revelaram três falhas estruturais
do Classifier Guidance nesse cenário (Seção~\ref{sec:cg-falha}), o que
motivou a migração para uma formulação inteiramente baseada em CFG
composicional~\cite{liu2022}, na qual identidade e atributos são
condicionantes independentes com escalas de guidance próprias. Este
relatório documenta a arquitetura, o protocolo de treinamento e a
sequência de experimentos que levou ao método final, implementado em
\texttt{train/train\_cfg\_composable\_paired.py}.

% =========================================================
\section{Dados}
% =========================================================

Todos os experimentos usam o \textbf{CelebA}~\cite{liu2015}: 202.599
fotos de rosto já alinhadas na convenção nativa do dataset, cada uma
anotada com 40 atributos binários (\emph{Smiling}, \emph{Eyeglasses},
\emph{Bald}, \ldots) e com um identificador de pessoa
(\texttt{identity\_CelebA.txt}). As identidades têm em média ${\sim}20$
fotos cada, e 91\% das identidades com pelo menos duas fotos variam no
atributo \emph{Smiling} --- propriedade que fundamenta o pareamento de
referência da Seção~\ref{sec:pareamento}.

O pré-processamento segue uma convenção única em todo o pipeline:
\texttt{CenterCrop(178)} seguido de \texttt{Resize(256)} sobre o recorte
nativo do CelebA. Nenhum script de treino de difusão lê imagem crua; três
caches pré-computados alimentam o treinamento:

\begin{itemize}
  \item \textbf{Cache de latentes} --- saída determinística do encoder do
    VAE (média $\mu$, sem amostragem), escalada por $0{,}18215$, com
    formato $[4, 32, 32]$ por imagem;
  \item \textbf{Cache de encoder} --- embedding ArcFace $[512]$ e tokens
    de patch do CLIP ViT-B/32 $[49, 768]$ por imagem, evitando rodar os
    backbones congelados a cada época;
  \item \textbf{Cache alinhado do classificador} --- recortes
    $224{\times}224$ usados apenas pelo classificador de atributos
    (Seção~\ref{sec:classificador}), com convenção de alinhamento
    própria.
\end{itemize}

Para o classificador de atributos, o split treino/validação/teste é feito
\emph{por identidade}, não por imagem, evitando que a mesma pessoa vaze
entre conjuntos. O CelebA é fortemente desbalanceado em vários atributos
(p.\,ex.\ \emph{Bald} tem ${\sim}2\%$ de positivos), o que motivou o uso
de Asymmetric Loss~\cite{benbaruch2021} no treino do classificador.

% =========================================================
\section{Arquitetura}
% =========================================================

\subsection{Visão geral}

O sistema é um LDM condicional com três componentes treináveis --- a
U-Net de denoising, o projetor de condicionamento de imagem e o embedder
de atributos --- sobre backbones congelados (VAE, ArcFace, CLIP). O
contexto de condicionamento é uma sequência de tokens de dimensão 512,
consumida pela U-Net via cross-attention; nenhuma modificação na U-Net é
necessária para adicionar ou remover ramos de condicionamento, pois o
cross-attention atende a todos os tokens concatenados. A
Tabela~\ref{tab:componentes} resume os componentes.

\begin{table*}[t]
  \centering
  \caption{Componentes do modelo. Parâmetros congelados não recebem
  gradiente.}
  \label{tab:componentes}
  \small
  \begin{tabular}{llrl}
    \toprule
    \textbf{Componente} & \textbf{Configuração} & \textbf{Parâmetros} & \textbf{Treinável} \\
    \midrule
    VAE                           & latente $4{\times}32{\times}32$, fator $0{,}18215$            & 2,5\,M  & não (pré-treinado) \\
    U-Net condicional             & base 128, mults $(1,2,4)$, atenção em $16^2$ e $8^2$          & 87,7\,M & sim \\
    ArcFace (buffalo\_l)          & ONNX, embedding $[512]$                                       & ---     & não \\
    CLIP ViT-B/32                 & tokens de patch $[49, 768]$                                   & ${\sim}88$\,M & não \\
    Projeções de condicionamento  & ArcFace $\to$ 16 tokens; CLIP $\to$ 16 tokens                 & 4,9\,M  & sim \\
    Embedder de atributos         & 40 tokens aprendidos $+$ \texttt{Embedding(2, 512)}           & 21,5\,K & sim \\
    \bottomrule
  \end{tabular}
\end{table*}

\subsection{U-Net de denoising}

A U-Net opera no latente $4{\times}32{\times}32$ com canais base 128 e
multiplicadores $(1, 2, 4)$, dois blocos residuais por nível e embedding
temporal senoidal projetado por MLP. Nas resoluções $16{\times}16$ e
$8{\times}8$ (e no gargalo), cada bloco recebe um
\texttt{SpatialTransformer}: self-attention seguida de cross-attention
sobre o contexto de tokens e uma feed-forward $4{\times}$, com projeções
de saída inicializadas em zero para estabilidade no início do treino.

\subsection{Condicionamento}

Três fontes de condicionamento são projetadas para o mesmo espaço de
contexto (dim.~512):

\begin{itemize}
  \item \textbf{Identidade (ArcFace).} O embedding L2-normalizado $[512]$
    é projetado linearmente em 16 tokens de contexto. O backbone ArcFace
    é totalmente congelado (ONNX).
  \item \textbf{Aparência (CLIP).} Os 49 tokens de patch do CLIP ViT-B/32
    (sem o token CLS) são reduzidos a 16 tokens por average pooling
    adaptativo e projetados para dim.~512 --- estilo
    IP-Adapter~\cite{ye2023}. Esse ramo captura cabelo, iluminação, pose
    e fundo da referência.
  \item \textbf{Atributos.} Cada um dos 40 atributos binários vira um
    token: um vetor aprendido por atributo somado a um embedding do valor
    binário (\texttt{nn.Embedding(2, 512)}), produzindo contexto
    $[B, 512, 40]$.
\end{itemize}

O contexto final é a concatenação dos ramos na dimensão de tokens ---
72 tokens no modo padrão (16 CLIP $+$ 16 ArcFace $+$ 40 atributos).

\subsection{CFG composicional}
\label{sec:cfg-composicional}

Na inferência, a predição de ruído combina múltiplos forward passes da
U-Net segundo a formulação composicional de Liu et~al.~\cite{liu2022}:
\begin{equation}
\label{eq:composable}
\begin{split}
\hat{\epsilon} = \; & \epsilon(\varnothing,\varnothing)
  + s_{\mathrm{id}}\left[\epsilon(\mathrm{id},\varnothing)
      - \epsilon(\varnothing,\varnothing)\right] \\
  & + s_{\mathrm{attr}}\left[\epsilon(\mathrm{id},\mathrm{attr})
      - \epsilon(\mathrm{id},\varnothing)\right]
\end{split}
\end{equation}

A ordem dos termos importa: o termo de atributo usa
$\epsilon(\mathrm{id},\varnothing)$ como baseline --- mede ``dada a
identidade, quanto o atributo desloca a predição'' --- em vez de competir
com o termo de identidade. As escalas $s_{\mathrm{id}}$ e
$s_{\mathrm{attr}}$ são ajustáveis independentemente na inferência,
recuperando o controle separado que a formulação Mixed Guidance oferecia
via \texttt{cfg\_scale} e \texttt{cg\_scale}.

No modo \textbf{split} (três ramos), ArcFace e CLIP deixam de ser
fundidos em um único condicionante ``id'':
\begin{equation}
\label{eq:split}
\hat{\epsilon} = \epsilon(\varnothing,\varnothing,\varnothing)
  + s_{\mathrm{id}}\,\Delta_{\mathrm{id}}
  + s_{\mathrm{clip}}\,\Delta_{\mathrm{clip}}
  + s_{\mathrm{attr}}\,\Delta_{\mathrm{attr}}
\end{equation}
onde cada $\Delta$ é a diferença de predições encadeada (identidade
primeiro, depois aparência, depois atributos). Zerar $s_{\mathrm{clip}}$
remove inteiramente a aparência da referência mantendo identidade e
controle de atributo --- essencial em edições estruturais como
\emph{Bald}, nas quais o cabelo visível na referência conflita
diretamente com o atributo desejado.

% =========================================================
\section{Treinamento}
% =========================================================

\subsection{Objetivo e difusão}

O processo direto é um DDPM padrão~\cite{ho2020} com 1.000 passos e
schedule linear de $\beta$ ($10^{-4}$ a $2{\cdot}10^{-2}$). O modelo é
treinado com o objetivo simples de predição de ruído sobre latentes:
\begin{equation}
\label{eq:loss}
\mathcal{L} = \mathbb{E}_{z_0,\,\epsilon,\,t}
  \left\| \epsilon - \epsilon_\theta(z_t, t, c) \right\|^2
\end{equation}
com $t$ uniforme em $[1, 1000)$ e contexto $c$ montado como na
Seção~\ref{sec:cfg-composicional}. O VAE permanece congelado durante todo
o treino.

\subsection{Pareamento de referência por identidade}
\label{sec:pareamento}

Um achado central do projeto: quando a imagem de referência do
condicionamento é a \emph{mesma} foto que é alvo do denoising, o ramo
CLIP enxerga a expressão exata que o modelo precisa reconstruir. O modelo
aprende a copiar a expressão da referência e ignora o token de atributo
--- na prática, nenhum valor de $s_{\mathrm{attr}}$ conseguia fazer os
rostos sorrirem. A correção é amostrar a referência de uma foto
\textbf{diferente da mesma identidade} sempre que a pessoa tem ${\geq}2$
fotos, mantendo o vetor de atributos da loss vindo do alvo. Como 91\% das
identidades elegíveis variam em \emph{Smiling}, a maioria dos batches
contém pares com expressões distintas, forçando o modelo a ler a
expressão do token de atributo e não da referência.

\subsection{Dropout de condicionamento}
\label{sec:dropout}

Para que a fórmula composicional disponha de todas as marginais na
inferência, o dropout de CFG é sorteado \emph{por amostra} (não por
batch): 10\% zeram só a identidade, 10\% zeram só os atributos, 10\%
zeram ambos e 70\% mantêm tudo. No modo split, cada um dos três ramos
recebe um dropout Bernoulli independente de 10\%. Zerar um ramo significa
anular seus tokens de contexto.

\subsection{Otimização}

A Tabela~\ref{tab:hiperparametros} resume os hiperparâmetros. Amostras
qualitativas e FID usam sempre os pesos EMA. A amostragem de avaliação
substitui os ${\sim}999$ passos ancestrais por DDIM~\cite{song2021} com
50 passos determinísticos --- cada passo custa 3 forwards da U-Net no
modo padrão (Eq.~\ref{eq:composable}) ou 4 no modo split
(Eq.~\ref{eq:split}). O último salto do DDIM vai direto para a imagem
limpa usando a predição de $x_0$ do último passo real, já que o modelo
nunca é treinado em $t = 0$.

\begin{table}[t]
  \centering
  \caption{Hiperparâmetros de treinamento
  (\texttt{train\_cfg\_composable\_paired.py}).}
  \label{tab:hiperparametros}
  \small
  \begin{tabular}{@{}p{0.42\linewidth}p{0.5\linewidth}@{}}
    \toprule
    \textbf{Hiperparâmetro} & \textbf{Valor} \\
    \midrule
    Otimizador          & AdamW, lr $3{\cdot}10^{-4}$, weight decay $10^{-4}$ \\
    Schedule de LR      & warmup linear (10 épocas) $+$ cosine annealing até $10^{-6}$ \\
    Batch / acumulação  & 64 por GPU, 4 passos de acumulação de gradiente \\
    Precisão            & mista (AMP) com \texttt{GradScaler} \\
    Paralelismo         & DDP multi-GPU (NCCL), \texttt{no\_sync} durante acumulação \\
    EMA                 & decay $0{,}9999$ sobre U-Net, encoder e embedder \\
    Dropout de CFG      & $10\%/10\%/10\%$ por amostra (Seção~\ref{sec:dropout}) \\
    Escalas de validação & $s_{\mathrm{id}}{=}3{,}0$; $s_{\mathrm{attr}}{=}5{,}0$; $s_{\mathrm{clip}}{=}3{,}0$ (split) \\
    Amostragem de avaliação & DDIM 50 passos, $\eta = 0$ \\
    \bottomrule
  \end{tabular}
\end{table}

\subsection{Protocolo de FID}

O FID~\cite{heusel2017} é calculado periodicamente durante o treino: as
estatísticas do lado real usam o conjunto de validação inteiro (extraídas
uma única vez e cacheadas em disco); o lado gerado usa 256 amostras
produzidas pelo modelo EMA a cada avaliação. O cálculo só começa após uma
época configurável (\texttt{-{}-fid\_start\_epoch}), pois nas primeiras
épocas o modelo ainda gera rostos incoerentes e o FID não é informativo.

% =========================================================
\section{Resultados e Experimentos}
% =========================================================

\subsection{Por que o Classifier Guidance falhou}
\label{sec:cg-falha}

A formulação original (``Mixed Guidance'') usava CFG para identidade e CG
para atributos, com um classificador
(\texttt{NoisyLatentAttrClassifier}) treinado sobre latentes ruidosos
$z_t$. Três problemas estruturais enfraqueceram o sinal de edição:

\begin{enumerate}
  \item \textbf{Saturação da sigmoide.} No classificador binário
    multi-rótulo, quando a predição está confiante (p.\,ex.\
    $\mathit{not\_smiling} \approx 1$) o gradiente de
    $\log p(y \mid z_t)$ colapsa; aumentar a escala de guidance para
    compensar gera artefatos antes de gerar o sorriso.
  \item \textbf{Diluição entre 40 atributos.} O gradiente soma os 40
    atributos --- os 39 não editados dominam e o sinal do atributo-alvo
    se dilui.
  \item \textbf{Janela temporal estreita.} Para $t$ alto o latente é
    quase ruído puro e o classificador aprende o prior; para $t$ baixo
    ele funciona, mas restam poucos passos para editar a imagem.
\end{enumerate}

Esses são precisamente os problemas que motivaram o
CFG~\cite{ho2022cfg}: a diferença de predições de ruído não satura e
escala de forma estável. A migração para CFG composicional eliminou o
classificador auxiliar, seu otimizador e a loss BCE do treino de difusão.

\subsection{Vazamento de expressão pelo CLIP}

Com CFG composicional e encoder CLIP$+$ArcFace condicionado na própria
imagem-alvo, o teste qualitativo padrão (forçar \emph{Smiling}=1 nas
amostras de validação e comparar com os atributos originais) revelou que
subir $s_{\mathrm{attr}}$ não fazia quase ninguém sorrir: o CLIP via a
expressão do alvo e o modelo a copiava. O pareamento por identidade
(Seção~\ref{sec:pareamento}) ataca a causa em vez do sintoma. Os
checkpoints permanecem intercambiáveis entre as variantes com e sem
pareamento, permitindo warm-start entre experimentos.

\subsection{Modo split para edições estruturais}

Mesmo com pareamento, o ramo CLIP ainda transporta aparência da
referência (cabelo, acessórios) que conflita com edições estruturais como
\emph{Bald}. O modo split (Eq.~\ref{eq:split}) separa ArcFace e CLIP em
ramos de guidance independentes usando os \emph{mesmos pesos} do encoder
conjunto --- na inferência, $s_{\mathrm{clip}} \to 0$ remove a aparência
da referência preservando identidade (ArcFace) e obediência ao atributo,
enquanto $s_{\mathrm{clip}}$ alto recupera reconstrução fiel de
iluminação e fundo. O ajuste vira uma escolha por edição, sem retreinar.

\subsection{Classificador de atributos}
\label{sec:classificador}

O classificador que fecha o pipeline (predizer os atributos da foto de
entrada para o usuário editar) é um CLIP ViT-L/14 com fine-tuning parcial
(últimos blocos $+$ head), treinado com Asymmetric
Loss~\cite{benbaruch2021}, warmup$+$cosine, clipping de gradiente e EMA,
sobre split por identidade. A variante com alinhamento pelo template
ArcFace atinge \textbf{F1\,$\approx$\,0{,}73} na validação, com limiar
por atributo calibrado para F1 ótimo; uma variante com o alinhamento
nativo do CelebA (que preserva a região de cabelo/chapéu que o recorte
ArcFace corta) está em comparação controlada.

\subsection{Métricas quantitativas e qualitativas}

A Tabela~\ref{tab:resultados} resume os resultados obtidos até o momento.
A cada checkpoint são salvas duas grades de amostras com as mesmas
identidades: uma com os atributos verdadeiros das referências e outra com
\emph{Smiling} forçado a 1 --- o contraste entre as duas é o teste direto
de que o modelo aprendeu a editar sem trocar a pessoa. O treinamento do
modelo final (\texttt{train\_cfg\_composable\_paired.py}, encoder
CLIP$+$ArcFace pareado) está em andamento; FID e taxa de acerto de
atributo verificada pelo classificador da Seção~\ref{sec:classificador}
serão consolidados ao término.

\begin{table}[t]
  \centering
  \caption{Resumo dos resultados obtidos até o momento.}
  \label{tab:resultados}
  \small
  \begin{tabular}{@{}p{0.5\linewidth}p{0.26\linewidth}r@{}}
    \toprule
    \textbf{Experimento} & \textbf{Métrica} & \textbf{Valor} \\
    \midrule
    DDPM incondicional $32{\times}32$ (baseline)      & FID                    & 3,35 \\
    Classificador de atributos (ViT-L/14)             & F1 macro (val.)        & ${\approx}\,0{,}73$ \\
    LDM CFG composicional pareado                     & FID (EMA, DDIM-50)     & em treino \\
    Edição \emph{Smiling}=1 vs.\ original             & qualitativa            & grades $4{\times}4$ \\
    \bottomrule
  \end{tabular}
\end{table}

% =========================================================
\section{Conclusão}
% =========================================================

Este projeto partiu da hipótese de que Classifier Guidance complementaria
o Classifier-Free Guidance na edição de atributos faciais e chegou, por
eliminação experimental, a uma conclusão diferente: para atributos
binários multi-rótulo sobre latentes ruidosos, o CG é estruturalmente
frágil, e o CFG composicional oferece o mesmo controle independente de
escalas com um sinal estável e sem classificador auxiliar. Duas lições
práticas se destacam. Primeiro, \emph{o que o encoder de condicionamento
enxerga durante o treino define o que o modelo obedece na inferência}:
condicionar na própria imagem-alvo ensina o modelo a ignorar o token de
atributo, e o pareamento por identidade corrige isso na raiz. Segundo,
separar os ramos de condicionamento (identidade pura vs.\ aparência) em
termos de guidance independentes transforma um conflito de treino em um
controle de inferência ajustável por edição.

Trabalhos futuros incluem consolidar as métricas quantitativas do modelo
pareado (FID e acurácia de atributo verificada pelo classificador),
comparar formalmente os encoders (CLIP$+$ArcFace vs.\ ArcFace-only vs.\
split) e explorar variantes arquiteturais já mapeadas --- dual
cross-attention no estilo IP-Adapter~\cite{ye2023} e modulação
adaLN/FiLM para atributos, que alcançaria todos os blocos da U-Net em vez
de apenas as resoluções com atenção.

% =========================================================
% Referências
% =========================================================
\begin{thebibliography}{12}
\small

\bibitem{ho2020}
J.~Ho, A.~Jain, P.~Abbeel.
\newblock Denoising Diffusion Probabilistic Models.
\newblock \emph{NeurIPS}, 2020.

\bibitem{liu2022}
N.~Liu, S.~Li, Y.~Du, A.~Torralba, J.~B. Tenenbaum.
\newblock Compositional Visual Generation with Composable Diffusion Models.
\newblock \emph{ECCV}, 2022.

\bibitem{rombach2022}
R.~Rombach, A.~Blattmann, D.~Lorenz, P.~Esser, B.~Ommer.
\newblock High-Resolution Image Synthesis with Latent Diffusion Models.
\newblock \emph{CVPR}, 2022.

\bibitem{ho2022cfg}
J.~Ho, T.~Salimans.
\newblock Classifier-Free Diffusion Guidance.
\newblock \emph{arXiv:2207.12598}, 2022.

\bibitem{dhariwal2021}
P.~Dhariwal, A.~Nichol.
\newblock Diffusion Models Beat GANs on Image Synthesis.
\newblock \emph{NeurIPS}, 2021.

\bibitem{deng2019}
J.~Deng, J.~Guo, N.~Xue, S.~Zafeiriou.
\newblock ArcFace: Additive Angular Margin Loss for Deep Face Recognition.
\newblock \emph{CVPR}, 2019.

\bibitem{radford2021}
A.~Radford et~al.
\newblock Learning Transferable Visual Models From Natural Language
  Supervision.
\newblock \emph{ICML}, 2021.

\bibitem{liu2015}
Z.~Liu, P.~Luo, X.~Wang, X.~Tang.
\newblock Deep Learning Face Attributes in the Wild.
\newblock \emph{ICCV}, 2015.

\bibitem{song2021}
J.~Song, C.~Meng, S.~Ermon.
\newblock Denoising Diffusion Implicit Models.
\newblock \emph{ICLR}, 2021.

\bibitem{benbaruch2021}
E.~Ben-Baruch et~al.
\newblock Asymmetric Loss for Multi-Label Classification.
\newblock \emph{ICCV}, 2021.

\bibitem{heusel2017}
M.~Heusel et~al.
\newblock GANs Trained by a Two Time-Scale Update Rule Converge to a Local
  Nash Equilibrium.
\newblock \emph{NeurIPS}, 2017.

\bibitem{ye2023}
H.~Ye, J.~Zhang, S.~Liu, X.~Han, W.~Yang.
\newblock IP-Adapter: Text Compatible Image Prompt Adapter for
  Text-to-Image Diffusion Models.
\newblock \emph{arXiv:2308.06721}, 2023.

\end{thebibliography}

\end{document}
